% Generated shared reviewed material.
% Full shared proofs; no book-only figure or new macro is required.
\section{Proofs for the finite-cutoff Gauss tail}
\label{gtl:details}
The arithmetic norm input is the explicitly pinned
\href{https://arxiv.org/abs/2610.04045v1}{de Faveri,
arXiv:2610.04045v1, Theorem~1.1 and Example~2.2}, at order three.
It is used on its actual primary cube-free ideal family.
We do not reprove that deep theorem. The primitive root conventions
are those already established in Proposition~\ref{cbk:root}, and
the Jacobi orientation is that of~\eqref{sw:character}.
The accompanying source map binds the accepted research proofs and
their independent reviews. Full third-party reading caches are not
public proof inputs.

\subsection{The coherent main retains all sixth-part and ray data}
\label{gtl:mean-proof}
For a good ideal $\mathfrak s$, the original fixed-source transport
in~\eqref{cbk:eta} gives
\begin{equation}\label{gtl:sixth-phase}
 \overline{\eta(\mathfrak s^6)}
   =\overline{\rho(\mathfrak s)}^{\,6}
                    \mathbf1_{(\mathfrak s,\mathfrak m\mathfrak n)=1}.
\end{equation}
This is an assertion about the total character on good arguments,
not about the order of the individual possibly angular/radial
$\xi_{\mathfrak m}$. Write $\psi=\bar\rho^{\,6}$, and let
$\delta_\psi$ be one if it is principal on good ideals and zero
otherwise. The conductor of $\psi$ belongs to fixed $S$ data.

Here is the fixed-ray counting input, with its full deletion.
For $t\ge0$, the classical disk estimate gives
\begin{equation}\label{gtl:ray-count}
 \begin{aligned}
 N_{\psi,q}(t)
 &:=\sum_{\substack{\mathfrak s\ {\rm good}\\
                    N\mathfrak s\le t,\;(\mathfrak s,\mathfrak q)=1}}
                                                  \psi(\mathfrak s)\\
 &=\delta_\psi\kappa_S g(\mathfrak q)t+E_{\psi,q}(t),\\
 |E_{\psi,q}(t)|
 &\ll_{S,\psi}t^{1/3}
                  \sum_{\mathfrak d\mid\mathfrak q}(N\mathfrak d)^{-1/3}
 \ll_\epsilon D^\epsilon t^{1/3},
 \qquad N\mathfrak q\le4D^2 .
 \end{aligned}
\end{equation}
We recall a proof to specify the uniformity being used.
For a fixed ideal modulus, count elements in each translated
lattice coset inside a disk of radius $R$. Convolve the disk
indicator with a fixed compact smooth nonnegative mollifier
of radius $\Delta$. The smoothed disks of radii $R-\Delta$
and $R+\Delta$ sandwich the original disk. Their zero modes
differ from its area by $O(R\Delta+\Delta^2)$.

For a nonzero dual vector $\xi$, polar integration gives
\begin{equation}\label{gtl:disk-fourier}
 \widehat{\mathbf1_{|z|\le R}}(\xi)
       =\frac{R J_1(2\pi R|\xi|)}{|\xi|},
 \qquad
 |J_1(u)|\ll\min\{1,u^{-1/2}\}.
\end{equation}
The decay follows directly from the angular integral for $J_1$:
for $u\ge1$, remove arcs of length $u^{-1/2}$ around the two
stationary points. Their total contribution is $O(u^{-1/2})$.
On the complement integrate once by parts in the phase;
the boundary terms and the integral of the derivative of
the reciprocal phase derivative are also $O(u^{-1/2})$.
For $u\le1$ the angular formula gives the bounded estimate.
Thus no unproved cancellation of lattice points is inserted.
The mollifier supplies arbitrarily rapid decay beyond
$|\xi|\asymp\Delta^{-1}$. Counting points of the fixed dual
lattice in unit-width radial shells yields
\begin{equation}\label{gtl:disk-error}
 \sum_{\xi\ne0}
    |\widehat{\mathbf1_{|z|\le R}}(\xi)\widehat\varphi(\Delta\xi)|
 \ll R^{1/2}\Delta^{-1/2}.
\end{equation}
The translation only changes unit phases.
Choose $\Delta$ comparable to $R^{-1/3}$ for large $R$.
The boundary and dual errors are both $O(R^{2/3})$.
Bounded $R$ is absorbed into the fixed constant.

A type-zero finite ray character is a fixed periodic generator
function, trivial on the six units. Apply the preceding count
to its finitely many residue cosets, and only then form the
possibly complex weighted sum. Division by six and the
covolume $\sqrt3/2$ of $\mathcal O$ give the coefficient
$\pi/(3\sqrt3)$. The principal coset average is the good-$S$
density; a nonprincipal character has mean zero.
For $t\ge1$ this proves the undeleted error $O(t^{1/3})$.
For $0<t<1$ the ideal count is zero and $t\le t^{1/3}$,
so the same ideal estimate remains valid; $t=0$ is exact.

Now use the \emph{entire} finite inclusion--exclusion:
\begin{equation}\label{gtl:full-IE}
 N_{\psi,q}(t)
  =\sum_{\mathfrak d\mid\mathfrak q}
       \mu(\mathfrak d)\psi(\mathfrak d)
                       N_{\psi,1}(t/N\mathfrak d).
\end{equation}
Divisors with $N\mathfrak d>t$ are still present, with their
subunit arguments. For principal $\psi$ the main sum is
$\kappa_St\prod_{p\mid q}(1-(Np)^{-1})$; otherwise it is zero.
The error in~\eqref{gtl:ray-count} follows. Its final divisor
product is $O_\epsilon(D^\epsilon)$ by separating finitely
many small primes and using $1+P^{-1/3}\le P^\eta$ at all
sufficiently large $P$. This is ideal-squarefreeness, not
rational-squarefreeness.

For fixed $\mathfrak h,\mathfrak r$, the original sum over
$V<N\mathfrak s\le b_{\mathfrak h\mathfrak r}$ has weight
$\dot G(c\,t^6)$. In the pure block the scale is
\begin{equation}\label{gtl:Stieltjes-scale}
 c=\frac{A_1N\mathfrak hN\mathfrak r}
              {(N\mathfrak k)^2N\mathfrak e F}.
\end{equation}
Closed prefixes at both endpoints give the exact error formula
\begin{equation}\label{gtl:Stieltjes-error}
 \begin{split}
 w(b)E_{\psi,q}(b)-w(V)E_{\psi,q}(V)
             -\int_V^b E_{\psi,q}(t)w'(t)\,dt,
 \qquad w(t)=\dot G(ct^6).
 \end{split}
\end{equation}
There are no half-weights at ideal norm ties. The angular
representations of $J_0,J_1$ imply, with
$\Lambda=A_1X/D^2$,
\begin{equation}\label{gtl:Bessel-variation}
 \|\dot G\|_\infty\ll1,\qquad
 |x\dot G'(x)|\ll\sqrt x,\qquad
 \int_V^b|w'(t)|\,dt\ll\sqrt\Lambda .
\end{equation}
Indeed $ct^6\le\Lambda$ and $b/V\le2$; differentiating
the Bessel factor, not $G$, proves the last assertion.
Consequently the error at each fixed outer set is
$O(D^\epsilon V^{1/3}(1+\sqrt\Lambda))$.

At the top window $N\mathfrak hN\mathfrak r<64$.
Every good ideal of norm below $64$ is sixth-free, because
a nontrivial good prime has norm at least seven. Therefore
the unique decomposition into its $S$ part and good part
runs exactly over all $\mathfrak w$ in~\eqref{gtl:actual-mean}.
The principal condition is precisely $\rho^6=1$.
On the original coprime support
$g(\mathfrak m\mathfrak n)=g(\mathfrak m)g(\mathfrak n)$.
This proves the stated finite formula for $M_{\rm top}/D$.

For the separate error, $|A_U|\le3^{\omega(\mathfrak m)}
\ll_\epsilon D^\epsilon$, and there are $O(D)$ original
annular ideals. After normalizing once by $D$,
$A_1/(D\sqrt F)\le A_1/D^2$ cancels the two column counts.
The $S$/ray/$\mathfrak e$ sets are fixed, the top
$\mathfrak w$ set is finite, and the \emph{original} finite
$\mathfrak k$ sum is majorized by $\sum(N\mathfrak k)^{-2}$.
Since $V^{1/3}\asymp X^{1/18}$, this proves
\eqref{gtl:ray-error}. Only its error uses these triangles;
the mean in~\eqref{gtl:actual-mean} stays signed.

\subsection{Exact good-modulus Gauss factors and separate primitive S values}
\label{gtl:root-proof}
The finite cubic sum at frequency $z$ is
\begin{equation}\label{gtl:frequency-sum}
 \frac1{\sqrt{N\mathfrak m}}\sum_{x\bmod\alpha_{\mathfrak m}}
       \chi_{3,\mathfrak m}(x)
        \exp\!\left(2\pi i\operatorname{Tr}
                   \frac{zx}{\delta\alpha_{\mathfrak m}}\right)
       =\overline{\chi_{3,\mathfrak m}(z)}G_{\mathfrak m}.
\end{equation}
For a unit frequency this is change of variable.
At a collision, one local additive frequency is zero and its
nonprincipal local cubic character has zero sum. CRT proves
the assertion for every good squarefree ideal, including
inert primes and nonsquarefree rational norms.
The normalized sum has modulus one at unit frequencies.

The total $\eta$ has type zero, so the previously proved
root formula is its finite Gauss sum divided by $\sqrt F$.
CRT at the coprime primitive conductors
$\mathfrak f_S,\mathfrak q$ gives exactly
\eqref{gtl:S-factor} times $\tau_q(\theta_q)/\sqrt{Nq}$.
Although their product is unit-trivial, the finite $S$ and
good-$q$ characters need not be unit-trivial separately.
No extra unit average is taken. The actual finite
$S$ character, including its conductor drop, is used at its
own primitive modulus. At unit conductor its Gauss sum is one.

On good squarefree $q$, all local powers
$\theta_q,\theta_q^2,\theta_q^3$ are nonprincipal and primitive.
The Gauss--Jacobi identity therefore gives
\begin{equation}\label{gtl:Gauss-Jacobi}
 \tau_q(\theta_q)\tau_q(\theta_q^2)
       =J_q(\theta_q,\theta_q^2)\tau_q(\theta_q^3).
\end{equation}
Over a residue field this follows by setting $x+y=t$ in the
product of Gauss sums; the zero-$t$ term vanishes because
the product character is nonprincipal. Multiplication by CRT
proves the composite identity: its additive cross multipliers
on the left combine to those of $\theta_q^3$ on the right.
The Jacobi sum itself has no additive-generator multiplier.

The correctly oriented result already proved in
Section~\ref{sw:core} is
\begin{equation}\label{gtl:Jacobi-identification}
 \begin{aligned}
 J_q(\theta_q,\theta_q^2)
     &=j(\mathfrak m)\overline{j(\mathfrak n)},\\
 j(\mathfrak m)/\sqrt{N\mathfrak m}
     &=\mu(\mathfrak m)\vartheta(\mathfrak m),
       \qquad \vartheta=\Psi/\sqrt N,\quad
                        \operatorname{cond}\Psi=(2\delta).
 \end{aligned}
\end{equation}
Its sign is fixed by $-j(p)\equiv1\pmod{2\delta}$ and
$(j(p))=p$ for the positive sextic character. The reduction
argument includes inert primes. Unique generators congruent
to one modulo $2\delta$ then identify $-j(p)$ with $\Psi(p)$.
Multiplication yields the displayed identity only on
squarefree columns. In particular $\Psi$ has its actual
infinity type $(1,0)$; it is not a finite-order gauge.

Cubic reciprocity for the same primary section and the fact
that $\chi_3(-1)=1$ yield
\begin{equation}\label{gtl:cubic-CRT-root}
 \begin{aligned}
 \chi_{3,\mathfrak m}(\alpha_{\mathfrak n})
   &=\chi_{3,\mathfrak n}(\alpha_{\mathfrak m}),\\
 \tau_q(\theta_q^2)/\sqrt{Nq}
   &=G_{\mathfrak m}\overline{G_{\mathfrak n}}.
 \end{aligned}
\end{equation}
The cubic cross factors cancel, while
$\tau(\bar\chi_3)=\overline{\tau(\chi_3)}$.
The quadratic Gaussian evaluation in the same two-factor
completion is
$\tau_q(\theta_q^3)/\sqrt{Nq}
 =\Gamma(\alpha_{\mathfrak m}\alpha_{\mathfrak n})$,
with the period-four formula in~\eqref{gtl:root-cubic}.
For composite $q$, this uses that Gauss evaluation with its
exact pairing; $1+\theta_q^3$ is not the root-count function
of a composite residue ring. Combining these facts with
\eqref{gtl:Gauss-Jacobi} proves~\eqref{gtl:root-cubic}.
There is no additional $\theta_q^3(-1)$ factor.

For the original good $\mathfrak k$,
\begin{equation}\label{gtl:frequency-transport}
 \begin{aligned}
 \eta(\mathfrak k^2)
   &=\theta_S(\alpha_{\mathfrak k})^2
        \chi_{3,\mathfrak m}(\alpha_{\mathfrak k})
        \overline{\chi_{3,\mathfrak n}(\alpha_{\mathfrak k})},\\
 \overline{G_{\mathfrak m}}G_{\mathfrak n}
        \chi_{3,\mathfrak m}(z)
        \overline{\chi_{3,\mathfrak n}(z)}
   &=\overline{g_{\mathfrak m}(z)}g_{\mathfrak n}(z).
 \end{aligned}
\end{equation}
Before the pure specialization the row factor
$R_{\mathfrak m}(\mathfrak v_2^2)
 \overline{R_{\mathfrak n}(\mathfrak v_2^2)}$
makes $z=\alpha_{\mathfrak k}\alpha_{\mathfrak v_2}$.
All raw collisions are zeros. The separate original
$(\mathfrak k,\mathfrak v_2)=1$ condition is not replaced by
a frequency assertion.
The primitive factors at $\mathfrak e$ and $\mathfrak w$ are
$\eta(\mathfrak e)$ and $\overline{\eta(\mathfrak w)}$.
When expressed with generators, their $S$ parts are the
separate factors $\theta_S(\gamma_{\mathfrak e})$ and
$\overline{\theta_S(\gamma_{\mathfrak w})}$.
If either is zero the component vanishes. No quotient of
these arguments, division by a primitive zero, or primary
section for an $S$-supported ideal is used. Changes of these
generators cancel in the total unit-trivial character.

\subsection{Matched finite endpoint models and a real-axis remainder}
\label{gtl:models-proof}
For $z\ge0$ and an integer $J\ge0$, the angular cosine formula
and its real Taylor remainder give
\begin{equation}\label{gtl:real-remainder}
 \left|J_0(2\sqrt z)
       -\sum_{j=0}^J\frac{(-z)^j}{(j!)^2}\right|
             \le\frac{z^{J+1}}{((J+1)!)^2}.
\end{equation}
Indeed use the cosine Taylor polynomial through degree
$2J+1$ inside its angular integral. The remainder is at most
$|2\sqrt z\cos\theta|^{2J+2}/(2J+2)!$, and the normalized
angular moment of $\cos^{2j}\theta$ is
$(2j)!/(4^j(j!)^2)$. This proves both the coefficients and
the bound for every nonnegative $z$, without an exponential
majorant or a sign assertion about the error.

Set
\begin{equation}\label{gtl:matched-moments}
 m_j=\int_0^\infty v^jG(v)\,dv,\qquad
 I_j(\mathfrak w)
   =\frac{b_{\mathfrak w}^{6j+1}-V^{6j+1}}{6j+1}
   =\int_V^{b_{\mathfrak w}}t^{6j}\,dt .
\end{equation}
The exact order-$J$ signed model of~\eqref{gtl:actual-mean} is
\begin{equation}\label{gtl:all-models}
 \mathcal L_J[f_1,f_2]
  =C_0^{-1/2}\sum_\theta d_\theta[f_1,f_2]
       \sum_{j=0}^J
        \frac{(-1)^jm_j\lambda_\theta^j I_j(\mathfrak w)}
             {C_0^j(j!)^2}\,
                   Q_{2+2j}^{T_q}(K;\eta^2).
\end{equation}
The character in every $Q_s$ is still $\eta^2$.
Both finite endpoints, every primitive root and every
original finite $\mathfrak k$ cutoff remain.
In particular $I_0=X^{1/6}t_{\mathfrak w}$ and
\eqref{gtl:all-models} at $J=0$ is exactly
\eqref{gtl:leading-model}.

Applying~\eqref{gtl:real-remainder} before the original
$\mathfrak k$ sum and the $t$ integral gives a per-$\theta$
bound
\begin{equation}\label{gtl:per-mode-error}
 \begin{split}
 \frac{C_0^{-1/2}m_{J+1}
        (\lambda_\theta/C_0)^{J+1}
        I_{J+1}(\mathfrak w)}
       {((J+1)!)^2}
     \sum_{\substack{\mathfrak k\ {\rm original}}}
                   (N\mathfrak k)^{-4-2J}.
 \end{split}
\end{equation}
Only this certified remainder takes absolute values.
On the full support
$\lambda_\theta b_{\mathfrak w}^6
 =A_1X/(N\mathfrak e F)\le\Lambda$.
Since $m_{J+1}\le2(5/2)^{J+1}$ and
$\sum(N\mathfrak k)^{-4-2J}$ is bounded by one constant
independent of $J$ and all deletions, put
$\Omega=(5/2)\Lambda/C_0$. Counting the two actual annular
column sets only in the error yields
\begin{equation}\label{gtl:all-model-error}
 |\mathcal M_{\rm top}-\mathcal L_J[A_U,A_U]|
 \ll B_W^2(DX)^\epsilon A_1X^{1/6}
             \frac{\Omega^{J+1}}{((J+1)!)^2}.
\end{equation}
Here the two $O(D)$ column counts cancel
$1/(D\sqrt F)\le D^{-2}$, the $S$/source/ray sets are fixed,
the $\mathfrak w$ set has bounded size, and each interval
has length at most $V$. The constant in this inequality is
independent of $J$ and of derivatives of $G$.
At $J=0$ it proves~\eqref{gtl:model-error}.
The infinite Bessel hierarchy was already available; the
statement used here is the finite real-axis certificate.
It bounds the remainder, not the signed polynomial.

\subsection{All source and ray phases give bounded multiplicative columns}
\label{gtl:phase-proof}
Freeze $\mathfrak n,\mathfrak e,\mathfrak w,\mathfrak k$
and one source/ray component in the pure model.
After~\eqref{gtl:root-cubic}, every $\mathfrak m$ dependence
apart from
$\mu(\mathfrak m)f(\mathfrak m)
 \overline{G_{\mathfrak m}}W(N\mathfrak m/D)/\sqrt{N\mathfrak m}$
is a product of bounded multiplicative functions, except
for two fixed finite selectors.
The multiplicative factors include $\vartheta,\nu,g$,
fixed-numerator residue symbols at $u_0,\gamma_S,\mathfrak k,
\mathfrak e,\mathfrak w$, the value
$\theta_S(\alpha_{\mathfrak m})$, and the original
$(\mathfrak m,\mathfrak n)=1$ deletion.
They include zeros: a numerator valuation divisible by six
does not restore a collision.
Primary generators multiply on the required coprime products.

Let $\mathcal C$ be the fixed source group and
$\mathcal H=(\mathcal O/(4))^\times$, of order $12$.
Finite character expansions give
\begin{equation}\label{gtl:finite-selectors}
 \begin{aligned}
 \mathbf1_{[\mathfrak m]=A}
   &=\frac1{|\mathcal C|}\sum_{\chi\in\widehat{\mathcal C}}
           \overline{\chi(A)}\chi([\mathfrak m]),\\
 \Gamma(x\alpha_{\mathfrak n})
   &=\sum_{\chi\in\widehat{\mathcal H}}a_{\chi,\mathfrak n}\chi(x),
       \qquad \sum_\chi|a_{\chi,\mathfrak n}|\le\sqrt{12}.
 \end{aligned}
\end{equation}
The first coefficient mass is one. The second follows by
normalized Parseval and Cauchy since the function has
modulus one on $\mathcal H$.
These bounds are independent of $\mathfrak n,\mathfrak k,D$.
Extending a fixed formula off its source class only extends
a coefficient function; the finite selector sum restores
the original class. No primitive root of a wrong-class
or overlapping-conductor pair is asserted.

Consequently every actual fixed-other-column expression is
a finite linear combination, of fixed total coefficient
mass, of the linear forms in~\eqref{gtl:linear-tail}.
The associated $\lambda$ may have growing conductor and
other-column deletion, but $|\lambda|\le1$ and
multiplicativity on coprime products suffice.
This does not flatten the original roots or make the
coprime matrix positive. It is not a silent extension of
a same-class, $\rho=1$ formula.

\subsection{Two exact cutoff boxes expose a physical cubic matrix}
\label{gtl:grouping-proof}
To prove~\eqref{gtl:ternary}, in each term of
\eqref{gtl:coefficients} set
$\mathfrak c=\mathfrak m/(\mathfrak a\mathfrak b)$.
Squarefreeness makes the factors pairwise coprime and
$\mu(\mathfrak m)\mu(\mathfrak a)\mu(\mathfrak b)
 =\mu(\mathfrak c)$.
The cutoff difference is the exact disjoint assignment
\begin{equation}\label{gtl:two-boxes}
 \begin{aligned}
 &\mathbf1_{N\mathfrak a,N\mathfrak b\le U}
       -\mathbf1_{N\mathfrak a,N\mathfrak b\le Y}\\
 &\quad=\mathbf1_{Y<N\mathfrak a\le U,\;N\mathfrak b\le U}
       +\mathbf1_{N\mathfrak a\le Y,\;Y<N\mathfrak b\le U}.
 \end{aligned}
\end{equation}
Swap $\mathfrak a,\mathfrak b$ in the second piece and
put $\mathfrak z=\mathfrak b\mathfrak c$ in both.
On coprime primary $\mathfrak a,\mathfrak z$, Gauss CRT
has the cross
$\chi_{3,\mathfrak a}(\alpha_{\mathfrak z})^2$.
Conjugation and cubic reciprocity give exactly
\eqref{gtl:gauss-cross}. Thus the left side of
\eqref{gtl:linear-tail}, before taking its absolute value,
is
\begin{equation}\label{gtl:grouped-form}
 \begin{split}
 \sum_{\substack{\mathfrak a,\mathfrak z\ {\rm good,SF}\\
                         Y<N\mathfrak a\le U}}
 &\frac{\overline{G_{\mathfrak a}}\lambda(\mathfrak a)}
             {\sqrt{N\mathfrak a}}\,
   \frac{\overline{G_{\mathfrak z}}\lambda(\mathfrak z)
                         [T_U(\mathfrak z)+T_Y(\mathfrak z)]}
             {\sqrt{N\mathfrak z}}\\
 &\quad\cdot\chi_{3,\mathfrak a}(\alpha_{\mathfrak z})
                         W(N\mathfrak aN\mathfrak z/D).
 \end{split}
\end{equation}
At a common prime the cubic entry is zero. Hence this
encompassing SF rectangle encodes the original
squarefreeness of $\mathfrak m=\mathfrak a\mathfrak z$,
without inventing a primitive overlapping root.
Also
$\mathbf1_{(\mathfrak a\mathfrak z,\mathfrak n)=1}
 =\mathbf1_{(\mathfrak a,\mathfrak n)=1}
  \mathbf1_{(\mathfrak z,\mathfrak n)=1}$, and the other
fixed-numerator collision masks split in the same way.
In particular the full other-column mask survives.
There is no assumption that $A_U$ or $T_L$ is multiplicative.
The elementary estimate
$|T_U(\mathfrak z)+T_Y(\mathfrak z)|\le2\tau(\mathfrak z)$
is uniform in both cutoffs.

\subsection{Cubic Cauchy before the distinguished-divisor triangle}
\label{gtl:cubic-proof}
Use dyads $N\mathfrak a\asymp A$ in $(Y,U]$.
Only $N\mathfrak z\asymp Z$, $Z\asymp D/A$,
can meet the original profile. Both indices in
\eqref{gtl:grouped-form} are good squarefree ideals, hence
cube-free. Primary cubic reciprocity identifies the entry
with the physical family of~\eqref{gtl:cubic-input}, or
its transpose. Additional row/column masks are coefficient
contractions inside that arbitrary-coefficient norm.
No character indexed only by a growing conductor is substituted.

Ideal counting and the divisor estimate show, with
arbitrarily small exponent losses,
\begin{equation}\label{gtl:dyadic-L2}
 \begin{aligned}
 \sum_{N\mathfrak a\asymp A}
   \left|\frac{\overline{G_{\mathfrak a}}\lambda(\mathfrak a)}
                    {\sqrt{N\mathfrak a}}\right|^2&\ll1,\\
 \sum_{N\mathfrak z\asymp Z}
   \left|\frac{\overline{G_{\mathfrak z}}\lambda(\mathfrak z)
                    [T_U(\mathfrak z)+T_Y(\mathfrak z)]}
                    {\sqrt{N\mathfrak z}}\right|^2
                                                    &\ll_\epsilon D^\epsilon .
 \end{aligned}
\end{equation}
These are coefficient bounds after the exact signed grouping,
not a new bound for the norm of $A_U$.

The exact Mellin separation of the one common occurrence
of $W$ is
\begin{equation}\label{gtl:Mellin-separation}
 W(N\mathfrak aN\mathfrak z/D)
  =\frac1{2\pi}\int_{\mathbb R}\widehat W(t)D^{-it}
                  (N\mathfrak a)^{it}(N\mathfrak z)^{it}\,dt .
\end{equation}
Its powers are unitary on this line, so it costs
$\mathcal M(W)$ once. The arbitrary-coefficient cubic
operator norm and~\eqref{gtl:dyadic-L2} therefore give
\begin{equation}\label{gtl:dyadic-column}
 |\text{one dyadic form}|
  \ll_\epsilon \mathcal M(W)D^\epsilon
                \sqrt{A+D/A+D^{2/3}} .
\end{equation}
This is a bilinear matrix inequality, not positivity of
the original coprime form.
Summing the finite dyads with $A\ge Y/2$ and
$A\le U\le\sqrt D$ gives
\begin{equation}\label{gtl:summed-column}
 \mathcal M(W)D^\epsilon
       \{\sqrt U+\sqrt{D/Y}+D^{1/3}\}
 \ll\mathcal M(W)D^\epsilon
       \{\sqrt{D/Y}+D^{1/3}\}.
\end{equation}
Dyadic logarithms and divisor losses are absorbed by
distributing a prescribed $\epsilon$ among the preceding
bounds. Partial dyads and real cutoff ties are literal
coefficient indicators. Empty tails, including $Y=U$,
are exactly zero. This proves~\eqref{gtl:linear-tail}.

\subsection{Full pair support and the original finite k reassembly}
\label{gtl:reassembly-proof}
Fix the other column $\mathfrak n$. Its raw coprimality
and numerator deletions are already in $\lambda$ by
Section~\ref{gtl:phase-proof}.
An optional condition $N\mathfrak m\ne N\mathfrak n$ is
not multiplicative. The number of ideals of a fixed
rational norm $N$ is at most the integer divisor function:
at a split prime an exponent $e$ gives at most $e+1$
choices, at an inert prime at most one, and at a
ramified prime one. Consequently there are
$O_\epsilon(D^\epsilon)$ ideals in the forbidden norm
fibre, with all orientations separate.
Since $|A_U|+|A_Y|\ll_\epsilon D^\epsilon$, subtracting it
costs
\begin{equation}\label{gtl:equal-fibre-cost}
 O_\epsilon(B_WD^{-1/2+\epsilon})
 \quad\text{in the fixed-other-column sum}.
\end{equation}
This is below~\eqref{gtl:linear-tail};
$\mathcal M(W)\ge B_W$ by Fourier inversion.
No whole equal-norm energy is subtracted from a different
frequency sector.

The untouched original column has the elementary bound
\begin{equation}\label{gtl:other-column}
 \sum_{\mathfrak n\ {\rm original}}
       \frac{|A_T(\mathfrak n)W(N\mathfrak n/D)|}
            {\sqrt{N\mathfrak n}}
       \ll_\epsilon B_WD^{1/2+\epsilon},
                 \qquad T\in\{U,Y\}.
\end{equation}
Take this triangle only after~\eqref{gtl:linear-tail}.
In the leading model all $S$, source/ray and
$N\mathfrak w<64$ coefficient masses are fixed;
$g_0$ is bounded independently of capacity derivatives,
and $F^{-1/2}$ supplies the two square-root norm weights
already used. The normalization is once by $1/D$.
For each actual fixed $\mathfrak k$, the column estimate
is uniform in its bounded multiplicative numerator
phase. Thus its \emph{original finite} sum is bounded by
$\sum(N\mathfrak k)^{-2}<\infty$, with no
$(N\mathfrak k)^\epsilon$ or altered conductor cost.
One does not need to complete that sum.

Equivalently, define $\mathcal G_D^K$ by
$\mathcal L_0[f_1,f_2]=A_1X^{1/6}\mathcal G_D^K[f_1,f_2]$.
The preceding estimates give
\begin{equation}\label{gtl:one-sided-moment}
 |\mathcal G_D^K[A_U-A_Y,A_T]|
 \ll B_W\mathcal M(W)D^\epsilon
                    \{D^{-1/6}+Y^{-1/2}\}.
\end{equation}
For a tail on the second column, conjugate the scalar
expression and reverse the roles. This sends $G_{\mathfrak n}$
to $\overline{G_{\mathfrak n}}$, $W$ to $\bar W$ and
the finite coefficients to their conjugates.
The conjugate cubic matrix has the same norm, and
$\mathcal M(\bar W)=\mathcal M(W)$.
No Hermitian assertion about an individual source block
is required. Finally use the exact bilinear identity
\begin{equation}\label{gtl:bilinear-difference}
 \begin{aligned}
 \mathcal G_D^K[A_U,A_U]-\mathcal G_D^K[A_Y,A_Y]
  &=\mathcal G_D^K[A_U-A_Y,A_U]\\
  &\quad+\mathcal G_D^K[A_Y,A_U-A_Y].
 \end{aligned}
\end{equation}
This proves~\eqref{gtl:finite-tail} with all original
finite cutoffs, primitive $S$ zeros and pair masks.

\subsection{The critical reduction is not a bound for its residual mean}
\label{gtl:critical-proof}
Equations~\eqref{gtl:model-error} and
\eqref{gtl:finite-tail} prove~\eqref{gtl:combined}.
When $U\ge D^{1/8}$, substituting
$A_1\asymp D^{3/4}$, $X=D^{3/8}$, $Y=D^{1/8}$
gives the exact exponent calculations
\begin{equation}\label{gtl:critical-powers}
 \begin{aligned}
 2(3/4)+(7/6)(3/8)-2&=-1/16,\\
 3/4+(1/6)(3/8)-(1/2)(1/8)&=3/4,\\
 3/4+(1/6)(3/8)-1/6&=31/48 .
 \end{aligned}
\end{equation}
If $U<D^{1/8}$ choose $Y=U$; the coefficient difference
then vanishes identically, rather than replacing
$Y^{-1/2}$ by a smaller number. The small-cutoff moment
remains exactly the auxiliary expression with that $Y$.
Since both short divisors have norms at most $Y$,
their residual in~\eqref{gtl:long-residual} has norm at
least $D/Y^2$. All other-column and finite-frequency
couplings remain in its two-column sum.

The current separate discrepancy has critical power
\begin{equation}\label{gtl:discrepancy-power}
 3/4+(1/18)(3/8)=37/48>3/4 .
\end{equation}
This upper certificate is not a noncancellation theorem.
A bound for the small-cutoff leading model alone does
not cancel or improve it. Both it and the remaining
original sectors require additional signed arithmetic.


\subsection{Direct discrete frequencies require no principal-density approximation}
\label{gtl:direct-proof}
For each original dual ideal, its good valuations determine
uniquely $\mathfrak y=\mathfrak w\mathfrak s^6$ with good
$\mathfrak s$ and good part of $\mathfrak w$ sixth-free.
At $V=X^{1/6}/2$, the lower-open condition
$N\mathfrak s>V$ implies $N\mathfrak w<64$.
Conversely every ideal of norm below $64$ has sixth-free
good part, since the least good prime norm is seven.
This proves the exact support~\eqref{gtl:discrete-support}
with all its nonunit $S$ modes and possible overlaps.
No squarefreeness of $\mathfrak s$ or new coprimality with
$\mathfrak k$ is introduced.

Equation~\eqref{gtl:sixth-phase} holds for \emph{every}
original $\rho$. Retaining that factor and the actual
discrete $\mathfrak s$ family in~\eqref{cbk:exact-kernel}
gives~\eqref{gtl:direct-prefix}. The prefactor is exactly
$A_1/(D\sqrt F)$, including the once-only normalization.
All primitive factors at $\mathfrak e,\mathfrak w$ are
still separate, and a good collision remains zero
even when its exponent is a sixth power.

The pointwise real-axis remainder~\eqref{gtl:real-remainder}
at $J=0$ implies
\begin{equation}\label{gtl:direct-Bessel}
 |\dot G(x)-g_0|
    \le C_0^{-3/2}x\int_0^\infty vG(v)\,dv
    \le3C_0^{-3/2}x,\qquad x\ge0 .
\end{equation}
Every argument in~\eqref{gtl:direct-prefix} is at most
$\Lambda=A_1X/D^2$. Count only in this error:
there are $O(D^2)$ annular column pairs,
$A_1/(D\sqrt F)\le A_1/D^2$, the original coefficient
products are $O(B_W^2D^\epsilon)$, and the original
finite $\mathfrak k$ absolute budget is bounded by
$\sum(N\mathfrak k)^{-2}$.
The $\mathfrak e$/source/ray and $\mathfrak w$ sets have
fixed finite mass. Since $V\ge1$, the exact good
$\mathfrak s$ annulus contains $O(V)$ ideals.
This proves~\eqref{gtl:direct-model-error}, without a
ray count, interpolated endpoint, density or arbitrary
row-weight norm.

It remains to justify applying the column-tail theorem
to every original ray. The good-$q$ powers in the proof
of~\eqref{gtl:root-cubic} are local powers $1,2,3$ of
$\theta_q$; that proof never uses $\rho^6=1$.
For every original component the total $\eta$ is type
zero and its finite $\theta_S$ is fixed at its actual
primitive $S$ conductor. Hence the finite selector
expansion in~\eqref{gtl:finite-selectors} and its
coefficient mass are valid on all these components.
No new assertion of sixth order for an individual
$\xi$ or replacement of a primitive $S$ Gauss sum is made.

Now freeze $\mathfrak s$ as well as
$\mathfrak n,\mathfrak k,\mathfrak e,\mathfrak w$.
Relative to Section~\ref{gtl:phase-proof}, the only new
$\mathfrak m$ factor is
$\mathbf1_{(\mathfrak m,\mathfrak s)=1}$.
It is a bounded multiplicative ideal function and
splits exactly in the grouped column as
\begin{equation}\label{gtl:s-mask-split}
 \mathbf1_{(\mathfrak a\mathfrak z,\mathfrak s)=1}
   =\mathbf1_{(\mathfrak a,\mathfrak s)=1}
      \mathbf1_{(\mathfrak z,\mathfrak s)=1}.
\end{equation}
The $\mathfrak n/\mathfrak s$ mask is a scalar which may
be zero. The fixed $\bar\rho(\mathfrak s)^6$ phase is
another bounded scalar. All earlier full-union raw
deletions and separate primitive $S$ values remain.
Thus~\eqref{gtl:linear-tail} is uniform in
$\mathfrak s$ as well as in $\mathfrak n,\mathfrak k$,
with no $(N\mathfrak s)^\epsilon$ or
$(N\mathfrak k)^\epsilon$ cost. Their growing deleted
primes are coefficient diagonals in the same physical
cubic matrix.

There is no density in the new exact discrete model.
Use the underlying arbitrary-$\lambda$ theorem with
the factor $g(\mathfrak m)$ omitted, rather than dividing
the earlier density-weighted moment by $g(q)$.
If a $\mathfrak k$ domain has been written with
$q$ deletion, first express it on the common original
good-SF set $N\mathfrak k\le K$ with its exact raw
$\eta(\mathfrak k)^2$ zero. This permits a fixed
$\mathfrak k$ column argument, restores no prime and
adds no $(\mathfrak k,\mathfrak s)$ condition.

Apply the signed column estimate before the other
column triangle. The latter is~\eqref{gtl:other-column};
after once-only $1/D$ it yields
$B_W\mathcal M(W)D^\epsilon
 (Y^{-1/2}+D^{-1/6})$ for each fixed
$\mathfrak s,\mathfrak w,\mathfrak k,\mathfrak e$
and finite component, with the original $A_1g_0$ factor.
Sum the original finite $\mathfrak k$ norm weights and
the fixed $\mathfrak e,\mathfrak w$ components.
Only now use the $O(V)$ count of the exact
$\mathfrak s$ family. Reversed tails use conjugation,
as in Section~\ref{gtl:reassembly-proof}, and the
original unequal-norm selector uses the same separate
fixed-norm fibre bound~\eqref{gtl:equal-fibre-cost}.
The exact identity
\begin{equation}\label{gtl:direct-bilinearity}
 \begin{aligned}
 \mathcal L_X^K[A_U,A_U]-\mathcal L_X^K[A_Y,A_Y]
  &=\mathcal L_X^K[A_U-A_Y,A_U]\\
  &\quad+\mathcal L_X^K[A_Y,A_U-A_Y]
 \end{aligned}
\end{equation}
then proves~\eqref{gtl:direct-model-tail}.
Combining it with the one Bessel error proves
\eqref{gtl:direct-bound}.

This direct proof does not assert a smaller value for
the separate $37/48$ certificate in the continuous
mean route; it never introduces that split.
All of its nonprincipal-ray and exact-discrete effects
remain in $\mathcal L_X^K[A_Y,A_Y]$.
At $Y\le D^{1/8}$ the residual column norms are at least
$D/Y^2\ge D^{3/4}$, just as in~\eqref{gtl:long-residual}.
The new model is still unbounded at target size.
It is different from the principal-density
$\mathcal L_0[A_Y,A_Y]$.

A capacity-empty integer interval, or empty exact
$\mathcal T_X$, gives zero by definition.
Absence of actual good squarefree rows in a
\emph{nonempty} integer interval alone does not imply
that each dual prefix piece is zero; cancellation
with other frequencies can be required.
No full dual complement, other block, whole signed
assembly or determinant/zero-location obligation is
resolved by this direct small-cutoff reduction.

\paragraph{Why the completed shortcuts do not supply that arithmetic.}
\label{gtl:failed}
The inverse-different cubic Gauss frequency is incomplete:
expanding it leaves the actual weighted primary
$\mathfrak k$ set, not a complete residue system modulo
$\mathfrak m\mathfrak n$. The bounded number of
$\mathfrak w$ modes supplies no such completion.
Likewise, expanding both short divisors and taking
their absolute triangles before the cubic cross
does not give the proved estimate. Even an explicitly
hypothetical uniform inner exponent $\alpha<1$ would
produce the diagnostic
\begin{equation}\label{gtl:triangle-diagnostic}
 D^{\alpha-1/2}
 \left(\sum_{N\mathfrak a\le U}(N\mathfrak a)^{-\alpha}\right)^2
 \ll_\alpha D^{\alpha-1/2}U^{2(1-\alpha)} .
\end{equation}
At $U\asymp\sqrt D$ this is $D^{1/2}$.
It is a conditional failure of that method, not an
available smooth-coefficient Gauss theorem or an
impossibility result. The actual proof above retains
the physical cubic cross until its joint matrix norm.
It bounds one specified signed tail and leaves the
long residual-M\"obius, finite-$\mathfrak k$ small-short
moment as the next arithmetic problem. No whole signed,
global unsigned, determinant, zero-location, analytic
Lean or GRH conclusion is implied.
